\summarychapter{30}{The Rule of the Saints}{Puritanism}

\begin{inbrief}
Puritanism, England's second reformation (about 1560--1660), joined an inward experience of converting
grace to a public program: a church and a nation ordered by the Bible. The chapter follows it through
three periods---the attempt to purify Elizabeth's church, resistance to James I and Charles I, and
the Civil War and Cromwell's rule---and notes the Separatists, Pilgrims and first Baptists who broke
away.
\question{What was Puritanism originally, and is the popular picture of the joyless Puritan fair?}
\end{inbrief}

In 1630 more than four hundred emigrants gathered at Southampton to sail for the New World. John
Cotton preached their farewell sermon on 2 Samuel 7:10, ``I will appoint a place for my people
Israel, and will plant them.'' Like ancient Israel, he said, they were God's chosen people bound for
a land he had prepared, where they could labor undisturbed for his glory. That is Puritanism in
miniature: the Bible as guidebook for the conversion of the individual and the construction of
society; the church, pictured as a new Israel, as the agent of reform; and the nation ordered to
serve God's purpose in history.

\ssection{Puritanism: New Life and New World}
Later Christians have often renewed the Puritan experience of the soul, but no later movement has had
the vision or the chance to apply God's law to the life of a nation. Today ``puritan'' means a
self-righteous, repressed killjoy; H.~L.~Mencken called Puritanism ``the haunting fear that someone,
somewhere may be happy.'' Shelley thinks this unfair. The Puritans were no killjoys; their movement meant change and the hope of a new age. Its plans grew from the conviction that spiritual conversion is crucial. This
rebirth set the Puritan apart from the mass of mankind and gave him the privileges and duties of the
elect. The church may prepare for it and guide afterward, but the reception of grace is beyond the
church's control.

The movement passed through three periods. In Elizabeth's reign (1558--1603) its aim was a Church of England cleansed after Geneva's pattern. Under James I and Charles I (1603--1642) it opposed what the crown demanded for itself and suffered pressure to conform to a high-church Christianity. During the Civil War and
Cromwell's rule (1642--1660) it had its chance to shape the national church and failed through its
own divisions. Its inward side made it a taproot of later evangelical, born-again Christianity; its
stress on a disciplined ``nation under God'' helped form the American character.

It began with a new kind of preaching, aimed at the heart rather than the head, by a ``spiritual
brotherhood'' (William Haller's phrase) that included Greenham, Rogers, Chaderton and Dod. Many had
been exiles under Bloody Mary and had found in Geneva a center of international subversion, a
``Protestant Moscow'' sending out students on fire to overthrow Catholicism at home. John Knox
succeeded in Scotland in 1560, and the English hoped for the same. They returned under Elizabeth
convinced that the Bible gives specific guidance for personal life, church structure and society. The
explosive issue was who chooses ministers: the Puritans said the people, the queen said the crown,
through bishops. Elizabeth let these dogs bark so long as they did not bite. They could complain of
set prayers, clerical dress and the sign of the cross in baptism, and lecture on preaching and the
office of elder, but the church stayed in the hands of bishops and so of the crown.

\ssection{A People of the Book}
Puritan zeal was fed by the Geneva Bible, the work of English exiles in that city, among them Miles
Coverdale, translator of the first complete English Bible in 1535. Its numbered verses, clear prose,
improved scholarship, prologues and marginal notes made it the most widely used English Bible until
the King James Version of 1611 replaced it, and the Puritans took it to America. With it went a sense
of destiny drawn from another product of the exile, John Foxe's \emph{Book of Martyrs}. Foxe traced a
trail of suffering witnesses down the centuries to England and the reign of Mary, and so suggested that the English had been given a part of their own in God's saving purpose. Next to
the Bible no book did more to shape the English mind, and the Puritans carried it too across the
Atlantic.

So they came to think of themselves as God's new Israel. The key is the biblical idea of covenant.
The most basic is the covenant of grace, the spiritual contract binding true Christians to God. God
sovereignly elects, yet anyone may by personal faith in Christ be added to the covenanted company. By
grace they become his people and he their God. The bond obliged the saint to walk in all God's ways
as made known in his Word, and it opened Scripture as a source of strength and of the rugged
determination for which Puritans became famous.

\ssection{Conflicts with the King}
Elizabeth died without an heir in 1603. James VI of Scotland, whose mother was Mary Queen of Scots, took the English throne as James~I and so joined the two crowns. Hopes that he would bring Presbyterianism south were soon
dashed. He had struggled with Presbyterian ministers for years and held that a Scottish presbytery
``agrees with monarchy as well as God and the devil.'' At the Hampton Court Conference in 1604 he
rudely dismissed the Puritans' proposals, granting only a new translation of the Bible, the King
James Version. Ceremonies, Prayer Book and bishops would stay, and the Puritans would conform or be
driven from the land. To Parliament he preached the divine right of kings, who are ``God's
lieutenants upon earth.'' He dissolved it in 1611 and ruled ten years without it, so that Puritans
and champions of Parliament drew together against royal power.

Some would not wait. Soon after Hampton Court small groups began to worship as they believed the
Bible taught, apart from bishops and prayer books. This was Separatism. Two such congregations, at
Scrooby and Gainsborough in northern England, had by 1608 moved to Holland, the first to Leyden, the
second to Amsterdam. After ten years the Scrooby group under Pastor John Robinson saw their children
losing touch with England and even its language, yet did not want to return to English society. They
had heard of the Virginia colony planted in 1607. A small number went back to England, joined others
of like mind, and in September 1620 about a hundred sailed from Plymouth in the \emph{Mayflower}. In
November they sighted New England, and the Pilgrim Fathers landed at what is now Plymouth,
Massachusetts.

The Gainsborough group was helped in Amsterdam by Mennonites, with whom they compared confessions.
Their pastor John Smyth, a Cambridge graduate, found no infant baptism in his Greek New Testament
and concluded that churches should be constituted by confession of faith. In 1609 he baptized
himself and then some forty members by pouring, forming the first Baptist church. Whether Baptists
are to be understood simply as English Separatists or as Separatists who mingled with Anabaptists
remains, Shelley notes, an open question of their origins and identity.

\ssection{Attack upon the Lord's Anointed}
Most Puritans found separation too drastic. But hope in earthly rulers vanished in 1625 with Charles
I, who meant to practice the divine right his father had only theorized about and whose wife,
Henrietta Maria, was a French Catholic princess. His instrument was Archbishop William Laud, who
believed that God ordained bishops to govern the church. Laud's party restored stained glass, crosses
and crucifixes, raised the communion table and called it an altar, and enforced the Prayer Book.
Some Puritans turned Separatist; others crossed the ocean. Within ten years of Laud's appointment
there were twenty towns and churches and sixteen thousand people in Massachusetts Bay.

The breaking point came when Charles tried to impose the \emph{Book of Common Prayer}, which Milton
called ``the skeleton of a Mass-book,'' on the Presbyterian Scots. They bound themselves in a
national covenant and took up arms. To raise an army Charles had to call the Parliament he had
ignored for over a decade. It split into Royalists and a Parliamentary majority that agreed on
Puritan principles but was divided between Presbyterians and Independents (Congregationalists).
United in hatred of Laud, they had him tried and beheaded. When Charles moved against their leaders,
civil war began, the Royalist members left, and Parliament was free to reform the church. The
Westminster Assembly of Puritan theologians (1643--1649) wrote the Westminster Confession of Faith, meant to take the place of the Thirty-nine Articles, and a Larger and a Shorter Catechism, still used by many
Presbyterians and Congregationalists. In 1645 Parliament ordered presbyteries and elders throughout
the land, but the system was never fully in force, for Parliament never ruled the whole country, and
after the king's defeat power lay with the Army.

\ssection{The Rule of the Saints}
The war brought forward Oliver Cromwell (1599--1658), a country gentleman who proved a military
genius. His regiment, the Ironsides, was never beaten, partly because he gave his men discipline and
a sense of Christian mission. He rose to lead the New Model Army, twenty-one thousand men who saw the
war as a crusade and went into battle praying and singing hymns. By the end of 1646 Charles had
surrendered. For two years he played off the Scots, the Presbyterians of Parliament and the
Independents of the Army against one another and made a secret alliance with the Scots. War broke out
again in 1648. The Army won, purged the Presbyterians from the Commons, and the remaining Rump
Parliament set up a court to try the king. In January 1649 Charles was beheaded before the palace of
Whitehall.

It was a rash act and a sign of Puritanism's coming fall, for it gave the royalists a martyr. A
widely circulated picture showed the king kneeling by a Bible, his crown on the floor, a crown of
thorns in his hand and his eyes on a crown of glory. It turned a king executed by rebels into a
martyr persecuted by fanatics, and that is how many Englishmen remembered the Puritans. (A chart in
the chapter is headed ``The Reformation in England.'') For the moment the saints ruled. The Commons
abolished the House of Lords and proclaimed a republic, the Commonwealth. In 1653 the Army, still
distrusting Parliament, replaced it with the Protectorate, with Cromwell as Lord Protector, in effect
a military dictator. He sought a religious settlement by granting liberty to Presbyterians,
Independents, Baptists, Quakers, Levellers and others, but found the task impossible, and his last
years were full of disappointment. His death in 1658 ended the saints' rule in old England. By 1660 the country had welcomed back the king, and with him the bishops. On the far side of the ocean, though, saints in their tens of thousands were building the kingdom of God in the wilderness of
New England.

\begin{keyfacts}
\keypeople{John Cotton; John Foxe; Miles Coverdale; Elizabeth I; James I; John Robinson; John Smyth;
Charles I; Archbishop William Laud; Oliver Cromwell.}
\keydates{%
  \kd{1603}{James VI of Scotland becomes James I of England}
  \kd{1604}{Hampton Court Conference}
  \kd{1609}{John Smyth forms the first Baptist church in Amsterdam}
  \kd{1611}{King James Version of the Bible}
  \kd{1620}{The Pilgrims sail in the \emph{Mayflower} and land at Plymouth}
  \kd{1625}{Charles I becomes king of England}
  \kd{1630}{Great Puritan migration to Massachusetts Bay begins}
  \kd{1643--1649}{Westminster Assembly: Confession and Catechisms}
  \kd{1649}{Charles I executed; the Commonwealth proclaimed}
  \kd{1653}{Cromwell becomes Lord Protector}
  \kd{1658}{Death of Cromwell}
  \kd{1660}{Monarchy and bishops restored in England}}
\keyterms{Puritanism; conversion; the elect; covenant of grace; new Israel; Geneva Bible; \emph{Book
of Martyrs}; divine right of kings; Separatism; Pilgrims; Presbyterians and Independents; Westminster
Confession; New Model Army; Commonwealth; Protectorate.}
\end{keyfacts}
